%Motivationsschreiben für den Master Studiengang Advanced Information Technologies
%
%
%Sehr geehrtes studienkomitee des Masterstudienganges Advanced Information Technologies,
%
%Als junge Person mit großem Interesse an Informationstechnologien und deren Anwendung in der modernen Welt, möchte ich mich hiermit für den Masterstudiengang Advanced Information Technologies an Ihrer renommierten Institution bewerben. Beim Absolvieren meines Bachelorstudiums in Informatik an der [Name der Universität] habe ich fundierte Kenntnisse in den Bereichen Programmierung, Datenbanken, Netzwerke und Softwareentwicklung erworben und dabei 
%
%Als junge erwachsende Person bin ich motiviert den Studiengang Advanced Information Technologien zu absolvieren. 
%
%Beim Absolvieren des Bachelors habe ich realisiert, dass ein Studium nur einem nicht inhaltliche Tiefe bezüglich eines Themas%
%Sehr geehrtes Studienkomitee des Masterstudienganges Advanced Information Technologies,%
%
%Meinen Bachelorstudiengang habe ich nach Inhaltlichen Kriterien ausgewählt. Dabei war es mir wichtig Wissen bezüglich der Informatik, Elektrotechnik und Softwareentwicklung stärker zu festigen. Was ich nicht erwartet habe, ist der Erwerb von Kompetenzen bezüglich der Problemlösung. Durch den regelmäßigen austausch mit Professoren und anderen% Studenten, erhielt ich Kompetenzen bezüglich des Denkens und des Lösen von Problemen. Ich bin ebenfalls der Aufassung, dass eine Ausbildung oder eine Festanstellung in einem Unternehmen mir nicht diese Kompetenzen vermittelen könnten. Mithilfe dieser Kompetenzen bin ich in der Lage komplexe Probleme zu bewaltigen.
%
%Aufgrund der positiven Erfahrungen im Bachelorstudiengang, würde ich mich gerne für den Masterstudiengang Advanced Information Technologies bewerben, weil ich überzeugt bin, dass ich als Person durch den Masterstudiengang mich weiter entwickeln kann im Vergleich zu anderen Optionen wie dem Arbeiten oder der direkten Selbständigkeit nach dem %Bachelor. Gleichzeitig strebe ich als Person in ferner Zukunft Jobmöglichkeiten mit Personalverantwortung an und hoffe ebenfalls Selbständigkeit agieren zu können. 



%-----------------------------------------------------------------------------%

% Letter class
\documentclass[a4paper, 10pt]{letter}

% Name of sender
\name{Erlind Sejdiu}

% Signature of sender
\signature{Erlind Sejdiu}

% Address of sender
\address
{
    Marktstraße 3,\\
    Ilsfeld, 74360 \\
}

%-----------------------------------------------------------------------------%

\begin{document}

\begin{letter}
{
    Erlind Sejdiu, Bachelorabsolvent
}

% Opening statement
\opening{Sehr geehrte Damen und Herren,}

meinen Bachelorstudiengang habe ich gezielt nach meinen persönlichen Interessen gewählt. Dabei war es mir wichtig, fundiertes Wissen in der Informatik, Elektrotechnik und Softwareentwicklung zu erwerben und stetig auszubauen. Während meines Studiums habe ich diese Chance konsequent genutzt und meine Kenntnisse vertieft. Durch die gezielte Wahl anspruchsvoller Module konnte ich tiefe Einblicke in die Metaprogrammierung gewinnen und praktische Erfahrungen in der Synthese von FPGAs sammeln. Ergänzend dazu habe ich im Bereich Software Engineering gelernt, hochkomplexe Programme strukturiert zu konzeptionieren und zu implementieren.

Ein für mich entscheidender, wenn auch unerwarteter Aspekt meines Studiums war die Entwicklung meiner Problemlösungskompetenz. Durch den regelmäßigen Austausch mit Professoren und Kommilitonen habe ich gelernt, Herausforderungen analytisch zu durchdringen. Die Art und Weise, wie ich heute Aufgaben angehe oder anderen komplexe Thematiken erkläre, ist maßgeblich dadurch geprägt, dass ich mich im Studium intensiv mit unterschiedlichen Abstraktionsebenen auseinandergesetzt habe. Meiner Auffassung nach ist der Erwerb von Fachwissen zwar essenziell, doch das gemeinsame Diskutieren, Reflektieren und Argumentieren mit Kollegen ist noch wertvoller. Ich bin überzeugt, dass eine klassische Ausbildung oder eine direkte Festanstellung mir diese spezifischen Kompetenzen in dieser Tiefe nicht hätten vermitteln können.

Aufgrund dieser positiven Erfahrungen im Bachelor möchte ich mich nun für den Masterstudiengang Advanced Information Technologies bewerben. Ich bin sicher, dass mich dieser Master als Persönlichkeit weiterbringt, da er im Vergleich zum direkten Berufseinstieg oder einer sofortigen Selbstständigkeit einen tieferen wissenschaftlichen Rahmen bietet. Besonders die Wahlmöglichkeiten im Bereich der Technischen Informatik passen ideal zu meinen aktuellen Schwerpunkten.

Meine Qualifikation für diesen Studiengang wird durch verschiedene Praktika unterstrichen. Bei der Firma Vector Informatik GmbH fokussierte ich mich auf Software Defined Networking im Automotive-Bereich. In einem weiteren Projekt konzeptionierte und implementierte ich ein Benachrichtigungsmodul für ein Betriebsdatenerfassungssystem.

Mit dem Master möchte ich meine Problemlösungsfähigkeiten weiter schärfen. Langfristig strebe ich eine Position mit Personalverantwortung oder eine unternehmerische Selbstständigkeit an. Da ein Masterabschluss oft die Voraussetzung für Führungspositionen ist, sehe ich dies als notwendigen Schritt. Auch privat arbeite ich leidenschaftlich an Softwareprojekten, etwa an einem Organisationstool für Projektumsetzungen sowie einer Applikation zur Simulation und zum Backtesting von Trading-Strategien. Diese Projekte haben mir gezeigt, dass ich in der Lage bin, Softwarelösungen eigenständig zu realisieren. Mein Ziel ist es, den gesamten Prozess von der Idee bis zur Vermarktung zu professionalisieren und mithilfe des Masters die nötigen Kompetenzen zu erwerben, um solche Projekte künftig noch effizienter und erfolgreicher umzusetzen.


----------------------------------------------------------------------------------------


nach meinem erfolgreich absolvierten Bachelorstudium an der Hochschule Pforzheim ist es mein Ziel, meine Kenntnisse in den Bereichen Informatik und Elektrotechnik gezielt zu vertiefen. Die positiven Erfahrungen aus meinem Erststudium sowie die spezifische Ausrichtung des Masters Advanced Information Technologies bestärken mich darin, meine akademische Laufbahn an Ihrer Hochschule fortzusetzen.

Im Laufe meines Bachelors habe ich mich intensiv mit komplexen Themen wie der Metaprogrammierung und der Synthese von FPGAs auseinandergesetzt. Besonders im Modul Software Engineering habe ich gelernt, anspruchsvolle Programme nicht nur zu implementieren, sondern von Grund auf systematisch zu konzeptionieren. Ein entscheidender Aspekt meines Studiums war jedoch nicht nur der reine Wissenserwerb, sondern die Entwicklung einer fundierten Problemlösungskompetenz. Durch die Auseinandersetzung mit unterschiedlichen Abstraktionsebenen habe ich gelernt, Herausforderungen analytisch zu durchdringen und alternative Lösungswege kritisch zu reflektieren eine Fähigkeit, die ich im engen Austausch mit Professoren und Kommilitonen stetig weiterentwickeln konnte.

Meine praktischen Erfahrungen untermauern meine fachliche Qualifikation. Während meines Praktikums bei der Vector Informatik GmbH spezialisierte ich mich auf Software Defined Networking (SDN) im Automotive-Sektor. Zudem konnte ich in einem weiteren Projekt ein Benachrichtigungsmodul innerhalb eines Betriebsdatenerfassungssystems eigenständig konzeptionieren und umsetzen. Diese Einblicke in industrielle Standards haben mir gezeigt, wie wichtig die wissenschaftliche Tiefe des Masterstudiums für die Lösung realer, hochkomplexer Probleme ist.

Neben dem Studium verfolge ich leidenschaftlich eigene Softwareprojekte. So habe ich ein Organisationstool für Projektmanagement sowie eine Applikation zur Simulation und zum Backtesting von Trading-Strategien entwickelt. Diese Projekte haben mir bewiesen, dass ich in der Lage bin, von der ersten Idee bis zur technischen Realisierung eigenständig zu agieren. Mein Ziel ist es, diese Projekte künftig zu professionalisieren und zur Marktreife zu führen.

Der Master Advanced Information Technologies ist für mich der nächste logische Schritt, um zwei Ziele zu vereinen: die technische Spezialisierung in der technischen Informatik und die Vorbereitung auf künftige Führungsaufgaben. Ich bin überzeugt, dass die akademische Tiefe dieses Studiums die ideale Basis bietet, um komplexe Projekte effizienter umzusetzen und später Verantwortung in leitender Position oder in der Selbstständigkeit zu übernehmen.

Ich freue mich darauf, meine Fähigkeiten und meine Motivation in den Masterstudiengang an der Hochschule Pforzheim einzubringen.

Mit freundlichen Grüßen

\closing{Mit freundlichen Grüßen,\\ Erlind Sejdiu}

\end{letter}
\end{document}
%-----------------------------------------------------------------------------%